\chapter{Problem Definition}
\label{ch:problem}

The central problem this project addresses is the \textbf{operationalization gap} in AI ethics: the significant disconnect between high-level ethical principles and the low-level technical and procedural mechanisms required to enforce them in practice. While numerous frameworks define \textit{what} constitutes ethical AI, they provide limited guidance on \textit{how} to systematically evaluate, measure, and govern these principles within a software development lifecycle~\cite{mittelstadt2019principles,hagendorff2020ethics}. This gap manifests in three primary sub-problems:

\begin{enumerate}[leftmargin=2em]
  \item \textbf{Lack of Standardized, Quantifiable Metrics.} Ethical concepts such as fairness, transparency, and accountability are abstract and context-dependent. Without a standardized methodology for translating these concepts into quantifiable metrics, evaluations remain subjective, inconsistent, and difficult to compare across different projects or over time. Developers lack a clear \enquote{yardstick} to measure their progress or to demonstrate compliance to regulators and the public.

  \item \textbf{Unmanaged Stakeholder Conflict.} Different stakeholders, including developers, business owners, regulators and the communities affected by the AI system, hold different values, priorities and risk tolerances~\cite{freeman1984stakeholder}. Although Freeman's stakeholder theory~\cite{freeman1984stakeholder} and Mitchell et al.'s stakeholder salience model~\cite{mitchell1997stakeholder} originate in general management, their core constructs (stakeholder identification, competing interests, and salience-based prioritization) map directly to AI governance, where developers, regulators and affected communities hold structurally analogous roles with divergent risk tolerances. A developer might prioritize system robustness and security, a regulator might focus on legal compliance and accountability, and an affected community member might be most concerned with fairness and the right to redress. A decision that appears optimal from one perspective may be unacceptable from another. This inherent multi-stakeholder conflict is a central challenge in AI governance~\cite{mitchell1997stakeholder}.

  \item \textbf{Absence of Auditable and Reproducible Workflows.} Ethical assessments are often conducted in an ad-hoc manner, with results recorded in static documents that lack traceability to the underlying data, assumptions, and decision logic. This makes it nearly impossible to reproduce an evaluation, audit a decision process, or track the evolution of a system's ethical posture over time. For AI governance to be effective, it requires a level of rigor and transparency comparable to other critical software engineering processes like testing and security~\cite{amershi2019se4ml}.
\end{enumerate}

Therefore, the problem statement for this project is as follows:

\begin{quote}
To design and implement a computational framework that enables organizations to systematically and reproducibly evaluate AI systems against multiple ethical governance standards, while explicitly modeling and managing the conflicting priorities of diverse stakeholders.
\end{quote}

To the best of the author's knowledge, no existing tool simultaneously evaluates AI systems against multiple governance frameworks, detects conflicts between stakeholder priorities, or produces stakeholder-weighted ethical scores. MEDF is designed to fill this gap.

To achieve this, the project sets the following objectives:

\begin{itemize}[leftmargin=2em]
  \item To develop a unified, six-dimensional model of AI ethics that harmonizes principles from leading international governance frameworks.
  \item To implement a multi-criteria decision-making (MCDM) engine using established algorithms (TOPSIS, WSM) to generate quantifiable ethical scores.
  \item To design and implement a mechanism for capturing stakeholder priorities as configurable weightings.
  \item To develop a conflict detection module that identifies and quantifies disagreement between stakeholder groups.
  \item To implement a multi-objective optimization algorithm (NSGA-II) to generate Pareto-optimal consensus solutions that help mediate stakeholder conflicts.
  \item To validate the framework's utility and effectiveness through its application to a set of diverse, real-world AI case studies.
  \item To deliver a completely documented, open-source software artifact that is both a practical tool and a platform for future research.
\end{itemize}
